\section*{Chapter \ref{ch:g11:subduction-zones} --- Subduction Zones}

\begin{solution}{exo:g11:subduction-zones:1}
\omterm{def:g11:subduction-zones:subduction}{Subduction} is the sinking of a plate of \omterm{def:g11:ocean-ridges:oceanic-lithosphere}{oceanic lithosphere} beneath another
plate, back into the mantle; a \omterm{def:g11:subduction-zones:subduction}{subduction zone} is a convergent plate boundary
where it happens, marked by a trench. \omterm{def:g11:ocean-ridges:oceanic-lithosphere}{Oceanic lithosphere} sinks: once it is
old, it is cold, thick and denser than the \omterm{def:g11:lithosphere-asthenosphere:lithosphere}{asthenosphere} beneath it. The
continental crust is too light to sink.
\end{solution}

\begin{solution}{exo:g11:subduction-zones:2}
(1) For example the Tonga--Fiji region, the western Pacific from Japan to the
Marianas, Indonesia, and South America. (2) Always on the side of the
overriding plate, never on the ocean side. (3) Shallow: all the earthquakes of
the Mid-Atlantic Ridge on the map are shallower than \qty{70}{km}.
\end{solution}

\begin{solution}{exo:g11:subduction-zones:3}
$7031 \times 0.21 \approx 1480$ earthquakes deeper than \qty{70}{km};
$7031 \times 0.069 \approx 485$ deeper than \qty{300}{km}.
\end{solution}

\begin{solution}{exo:g11:subduction-zones:4}
\omterm{def:g11:subduction-zones:eclogite}{Blueschist} is a metamorphic rock made from the basalt of the \omterm{def:g11:ocean-ridges:oceanic-lithosphere}{oceanic crust}
buried to high pressure while staying relatively cold, coloured blue by
glaucophane; \omterm{def:g11:subduction-zones:eclogite}{eclogite} is the dense metamorphic rock that basalt and \omterm{def:g11:ocean-ridges:gabbro}{gabbro}
become at mantle depths, made of red garnet and green pyroxene. Both come from
the \omterm{def:g11:ocean-ridges:oceanic-lithosphere}{oceanic crust} and keep its chemical composition; what differs are the
minerals, formed in the solid state under the new pressure and temperature,
and the density.
\end{solution}

\begin{solution}{exo:g11:subduction-zones:5}
Ocean floor: old \omterm{def:g11:lithosphere-asthenosphere:lithosphere}{lithosphere} moving toward the trench. Trench: the plate bends
down and starts to sink. Forearc: sediments scraped off the sinking plate pile
up; beneath, the two plates slide past each other on the plate boundary,
where great earthquakes happen. Volcanic arc: a line of volcanoes above the
\omterm{def:g11:subduction-zones:slab}{slab}. Back-arc: behind some arcs, the overriding plate is stretched.
\end{solution}

\begin{solution}{exo:g11:subduction-zones:6}
$\tan\delta = (508 - 135)/(800 - 300) = 373/500 \approx 0.746$, so
$\delta \approx \ang{37}$.
\end{solution}

\begin{solution}{exo:g11:subduction-zones:7}
(1) $\tan\ang{45} = 1$: \qty{300}{km} deep. (2) At \qty{670}{km} from the
trench. (3) $670\sqrt{2} \approx \qty{948}{km}$ of \omterm{def:g11:subduction-zones:slab}{slab}.
\end{solution}

\begin{solution}{exo:g11:subduction-zones:8}
(1) $(74.6 + 70.0 + 65.5 + 63.1 + 52.6)/5 \approx \qty{65.2}{mm}$ a year;
$(47.6 + 31.0 + 29.2 + 22.8 + 19.8 + 15.2 + 14.9 + 10.6)/8 \approx
\qty{23.9}{mm}$ a year; ratio $65.2/23.9 \approx 2.7$. (2) Plates move
fastest when a \omterm{def:g11:subduction-zones:slab}{slab} hangs from them: the \omterm{def:g11:subduction-zones:slab-pull}{pull of the sinking slab} is the
strongest force driving plates, as Forsyth and Uyeda concluded.
\end{solution}

\begin{solution}{exo:g11:subduction-zones:9}
(1) \qty{83}{mm} is \qty{0.083}{m}: $50/0.083 \approx 600$ years and
$60/0.083 \approx 720$ years of plate motion. (2) From 869 to 2011 is 1142
years: of the same order. If the plate boundary stays locked between great
earthquakes, the slip of 2011 released several centuries of stored motion.
\end{solution}

\begin{solution}{exo:g11:subduction-zones:10}
(1) Japan: $6400 - 3200\,e^{-133/62.8} \approx \qty{6015}{m}$; Tonga:
$6400 - 3200\,e^{-99/62.8} \approx \qty{5740}{m}$; Andes:
$2500 + 350\sqrt{49} = \qty{4950}{m}$. (2) The floor reaching Japan: the
oldest, so the coldest and densest \omterm{def:g11:lithosphere-asthenosphere:lithosphere}{lithosphere}, and the deepest floor.
\end{solution}

\begin{solution}{exo:g11:subduction-zones:11}
(1) $3380 \times 1.06 \approx \qty{3583}{kg/m^3}$. (2)
$(3583 - 2900)/2900 \approx 0.24$: a rise of about \qty{24}{\percent}.
(3) Before, the crust is lighter than the mantle by
$3380 - 2900 = \qty{480}{kg/m^3}$; after, it is denser by about
\qty{203}{kg/m^3}. It no longer holds the \omterm{def:g11:subduction-zones:slab}{slab} up: it pulls it down.
\end{solution}

\begin{solution}{exo:g11:subduction-zones:12}
(1) $t = 6 \times 10^6 \times 3.16 \times 10^7 \approx 1.9 \times 10^{14}$~s;
$\sqrt{10^{-6} \times 1.9 \times 10^{14}} \approx \qty{1.4e4}{m}$, about
\qty{14}{km}. (2) Much less than \qty{100}{km}: in 6 million years only the
outer skin of the \omterm{def:g11:subduction-zones:slab}{slab} warms up; its interior stays cold. (3) Cold rock is
rigid and can break, while the hot mantle around it flows; the deep
earthquakes happen inside the cold \omterm{def:g11:subduction-zones:slab}{slab}.
\end{solution}

\begin{solution}{exo:g11:subduction-zones:13}
(1) $\sqrt{1017^2 + 670^2} \approx \qty{1218}{km}$. (2) \qty{64}{mm} a year
is \qty{64}{km} per million years: $1218/64 \approx 19$ million years.
(3) About three times longer than beneath Tonga: the \omterm{def:g11:subduction-zones:slab}{slab} beneath the Andes
sinks more slowly and at a gentler dip.
\end{solution}

\begin{solution}{exo:g11:subduction-zones:14}
(1) A slice \qty{1}{m} wide along the trench has a volume of
$10^5 \times 10^6 \times 1 = \qty{1e11}{m^3}$; its excess mass is
$36 \times 10^{11} = \qty{3.6e12}{kg}$. (2) Its excess weight is
$3.6 \times 10^{12} \times 9.8 \approx \qty{3.5e13}{N}$ for each metre of
trench. (3) The \omterm{def:g11:subduction-zones:slab}{slab} is attached to the rest of the plate, which is rigid:
the force is transmitted to the whole plate, which is dragged toward the
trench, all the more since trenches run for thousands of kilometres.
\end{solution}

\begin{solution}{exo:g11:subduction-zones:15}
(1) $(23.6 - 16.5) \times 111 \approx \qty{790}{km}$. (2) The rate grows by
$237 - 104 = \qty{133}{mm}$ a year over \qty{790}{km}, about \qty{17}{mm} a
year per hundred kilometres. (3) The convergence is the velocity of the
Pacific plate minus that of the arc. Since the Pacific plate's velocity is
about the same at both points, the arc itself must be moving toward the
Pacific plate, faster and faster northward: the back-arc behind it is being
stretched, more in the north than in the south.
\end{solution}

\begin{solution}{pb:g11:subduction-zones:1}
\textbf{1.} About \qty{671}{km}.

\textbf{2.} On the ocean side the plate has not yet sunk: the \omterm{def:g11:lithosphere-asthenosphere:lithosphere}{lithosphere}
there is near the surface. The deep earthquakes happen inside the \omterm{def:g11:subduction-zones:slab}{slab}, which
lies only beneath the overriding plate.

\textbf{3.} $z(300) = -24 + 273 = \qty{249}{km}$ and
$z(700) = -24 + 637 = \qty{613}{km}$; $\tan\delta = 364/400 = 0.91$, so
$\delta \approx \ang{42}$.

\textbf{4.} $x = (100 + 24)/0.91 \approx \qty{136}{km}$ and
$x = (400 + 24)/0.91 \approx \qty{466}{km}$.

\textbf{5.} A straight line with this dip starting at the trench would pass
through $z = 0$ at $x = 0$; it lies above. The \omterm{def:g11:subduction-zones:slab}{slab} is not straight: near the
trench it dips more gently, and it steepens deeper down.

\textbf{6.} \qty{16.8}{cm} a year, \qty{168}{km} per million years.

\textbf{7.} $168/69 \approx 2.4$ times faster.

\textbf{8.} The convergence across the trench is the velocity of the Pacific
plate minus the velocity of the arc on the other side. If the arc moves toward
the Pacific plate, as it does while the back-arc behind it is stretched, the
two speeds add up, and the convergence is larger than the speed of the
Pacific plate alone.

\textbf{9.} $6400 - 3200\,e^{-99/62.8} \approx \qty{5740}{m}$.

\textbf{10.} After about 60 million years \omterm{def:g11:ocean-ridges:oceanic-lithosphere}{oceanic lithosphere} is denser, as a
whole, than the \omterm{def:g11:lithosphere-asthenosphere:lithosphere}{asthenosphere} beneath it; at 99 million years it is denser
still. Once its edge goes down, its weight pulls the rest in.

\textbf{11.} $(670 + 24)/0.91 \approx \qty{763}{km}$.

\textbf{12.} $\sqrt{763^2 + 670^2} \approx \qty{1015}{km}$.

\textbf{13.} $1015/168 \approx 6.0$ million years.

\textbf{14.} $1015/69 \approx 14.7$ million years.

\textbf{15.} $3380 \times 1.06 \approx \qty{3583}{kg/m^3}$: the crust, once
lighter than the mantle, becomes denser than it, and the whole slab pulls
harder.

\textbf{16.} $t \approx 6.0 \times 10^6 \times 3.16 \times 10^7 \approx
1.9 \times 10^{14}$~s, and $\sqrt{\kappa t} \approx \qty{14}{km}$.

\textbf{17.} \qty{14}{km} is small beside \qty{100}{km}: only the outer part
of the \omterm{def:g11:subduction-zones:slab}{slab} has warmed; its interior is still cold when it reaches
\qty{670}{km}.

\textbf{18.} The cold \omterm{def:g11:subduction-zones:slab}{slab} stays rigid enough to break, down to its tip,
while the hot mantle around it flows: the earthquakes beneath Tonga happen
inside the \omterm{def:g11:subduction-zones:slab}{slab}, down to about \qty{671}{km}.

\textbf{19.} Seismic waves show a sharp change in the mantle at about
\qty{670}{km}: their speeds jump there.

\textbf{20.} About \textbf{6 million years} from the trench to
\qty{670}{km}.
\end{solution}
